(1) 线性声学连续方程与Euler方程是
\[
\frac1{c^2}\partial_t p=-\frac{\rho}{S}\partial_x(Sv),\qquad
\rho\partial_t v=-\partial_x p.
\]
对第一式再求时间导数，代入第二式得
\[
p_{tt}={c^2\over S}\partial_x(Sp_x)=c^2(p_{xx}+2\alpha p_x).
\]
(2) 设 \(P=e^{-\alpha x}\psi\)，可得
\[
\psi''+k^2\psi=0,\qquad k^2={\omega^2\over c^2}-\alpha^2.
\]
传播分支 \(\omega>\alpha c\) 时 \(k\) 实且 \(\omega=c\sqrt{k^2+\alpha^2}\)；相速度 \(v_p=\omega/k\)，群速度 \(v_g=d\omega/dk=c^2k/\omega\)。截止 \(\omega_c=\alpha c\)。低于截止 \(k=i\kappa,\ \kappa=(\alpha^2-\omega^2/c^2)^{1/2}\)，被动向右衰减压力为 \(e^{-(\alpha+\kappa)x}\)，其平均净功率为零。
(3) 由 \(-i\omega\rho V=-P_x\)，对前行波
\[
V={k+i\alpha\over\rho\omega}P.
\]
平均声功率
\[
\overline{\mathcal P}=\frac S2\Re(PV^*)=\frac{S_0|A|^2k}{2\rho\omega}
\]
与位置无关。单位长度声能
\[
\overline{\mathcal E}=\frac S4\left({|P|^2\over\rho c^2}+\rho|V|^2\right)
=\frac{S_0|A|^2}{2\rho c^2},
\]
因为 \(k^2+\alpha^2=\omega^2/c^2\)。故 \(\overline{\mathcal P}/\overline{\mathcal E}=c^2k/\omega=v_g\)。当 \(\alpha\to0\) 时 \(k=\omega/c\)、\(V=P/(\rho c)\)、\(v_g=c\)。
(4) 写端面幅度 \(P(L)=e^{-\alpha L}(Ae^{ikL}+Be^{-ikL})\)，
\[
V(L)=\frac{e^{-\alpha L}}{\rho\omega}
\bigl[(k+i\alpha)Ae^{ikL}+(-k+i\alpha)Be^{-ikL}\bigr].
\]
定义 \(R_L=Be^{-ikL}/(Ae^{ikL})\)，代入 \(Z_L=P/V\) 后
\[
R_L=\frac{Z_L(k+i\alpha)-\rho\omega}{\rho\omega-Z_L(-k+i\alpha)}.
\]
匹配阻抗 \(Z_{\rm match}=\rho\omega/(k+i\alpha)\) 使 \(R_L=0\)；刚性端 \(V(L)=0\) 时 \(R_L=(k+i\alpha)/(k-i\alpha)\)，模长1。正反传播的净声功率与 \(|A|^2-|B|^2\) 成正比；若 \(\Re Z_L\ge0\)，端面吸收功率非负，故 \(|R_L|\le1\)。此结论对 \(\omega>\omega_c\) 的实波数成立，不能将倏逝模反射率作同一功率解释。
